\documentclass[conference]{IEEEtran}
\usepackage[utf8]{inputenc}
\usepackage[T1]{fontenc}
\usepackage{lmodern}
\usepackage{microtype}
\usepackage{hyperref}
\hypersetup{colorlinks=true,linkcolor=blue,urlcolor=blue,citecolor=blue}
\title{Loop Engineering: The Anthropic Playbook\\for Designing Systems That Prompt Your Agents}
\author{Reconstructed from publicly shared PDF text\\\textit{Not the original source TeX}}
\begin{document}
\maketitle
\begin{abstract}
This file is a best-effort, reconstructed LaTeX source generated from the publicly shared PDF text. It is provided for inspection and re-typesetting convenience only; it is not verified to be the original authoring source.
\end{abstract}
\section*{Extracted PDF Text}



% %%PAGE 1

Loop Engineering: The Anthropic Playbook
for Designing Systems That Prompt Your Agents
A Field Study of Designing Loops That Run Themselves
Abstract—Over the past two years a string of “XX Engineering”
terms has tracked the pace of model releases. This note examines
the newest of them, Loop Engineering, a term independently
surfaced in June 2026 by Peter Steinberger, Boris Cherny, and
Addy Osmani, and named in writing by Osmani. Unlike prompt,
context, or harness engineering, loop engineering does not teach
the practitioner to do the work better; it removes the practitioner
from the position of doing the work at all. We define the term,
place it as a fourth layer above the harness, and decompose a
single turn of a loop into five moves—discovery, hando ff, ver-
ification, persistence, and scheduling—and the six parts that
realize them. We give particular attention to the generator/eval-
uator separation: empirically, an agent asked to grade its own
output tends to praise it, and tuning an independent skeptical
evaluator is far more tractable than making a generator criti-
cal of its own work. We survey three loops running in practice,
from one engineer’s morning triage to Stripe’s enterprise-scale
pipeline merging over 1,300 machine-written pull requests per
week, and we catalog four costs that accrue silently—verification
debt, comprehension rot, cognitive surrender, and token blowout.
We close with a concrete recipe for building a first loop. The
central claim is that loops make generation nearly free and leave
judgment as the scarce resource; the same loop, built by two
people, can yield opposite outcomes.
Index Terms—Agentic AI, software engineering, autonomous
agents, coding agents, generator–evaluator, scheduling, automa-
tion.
I. I ntroduction: What Loop Engineering Really Is
The prompt engineering courses are still selling, the ink on
context engineering is not yet dry, and harness engineering has
only just been written up. Now loop engineering arrives as
well. Across the past year these “XX Engineering” terms have
appeared almost in step with model releases, and the temptation
to roll one’s eyes is understandable.
This one is different. It is not about doing the work better; it
is about pulling the practitioner out of the position of doing the
work entirely. The earlier terms all assumed a human seated at
the keyboard, directing the agent line by line. Loop engineering
deletes that assumption. The practitioner is no longer inside the
loop, but outside it, building the loop.
A. A One-Line Definition
The person who named the term and wrote it up is Addy
Osmani, an engineer on the Google Chrome team. His definition
is short: loop engineering is replacing oneself as the person
who prompts the agent, and designing the system that does it
instead. One no longer feeds the agent line by line; one designs
the system that feeds it. The weight of the sentence falls on
replacing yourself. This is a shift of position—from being the
engine to being the person who designs the engine. What one
writes is no longer the words for the agent, but a thing that
automatically sends words to the agent.
B. Within One Week, Three People Lit the Fuse
The term was not invented out of nowhere. In a single week
of June 2026, several groups ran into the same thing at almost
the same time. What set it off was a post from Peter Steinberger
(author of OpenClaw) which passed eight million views: one
should no longer be prompting coding agents, but designing
the loops that prompt them. At nearly the same moment Boris
Cherny (lead on Claude Code at Anthropic) was saying the same
thing—that he no longer prompts Claude, but has loops running
that prompt Claude and figure out what to do, and that his job
is to write loops. On June 7 Osmani wrote it up on his blog
under the title Loop Engineering, pulling in what Steinberger
and Cherny had said, and synced it to Substack the next day.
One ignition, one echo, one name, all within a week.
Stacked together, the three statements point at the same move:
what one designs has shifted from a single behavior of the agent
to the entire system that drives the agent. As with “vibe coding”
before it, the term may sound crude, but it turns a way of working
into something that can be discussed.
C. Why the Term Arrived Now
It is worth asking why three practitioners, not comparing
notes, reached for the same word in the same week. The answer
is that the surrounding tools had quietly crossed a threshold.
Coding agents had become reliable enough to finish a non-trivial
task unattended; scheduling primitives had appeared in the major
harnesses; and the cost of a single agent run had dropped far
enough that running one repeatedly, on a timer, stopped looking
wasteful. When the parts are all present, the move that combines
them becomes obvious to everyone at once. The name lagged the
practice by months: people were already writing loops before
anyone called it loop engineering, the same way teams were
already pairing a writer agent with a reviewer agent before the
generator/evaluator split had a name.
This pattern—practice first, name second—is worth keeping
in mind, because it tells the reader where to look for the next
term. It will not come from a model release. It will come
from the moment a new capability becomes cheap enough that a
previously unthinkable composition becomes routine.
D. One Floor Above the Harness
The shift reduces to two sentences. In the old world, one sits
and prompts the agent line by line; it finishes one thing, stops,
and waits. One is the human clock inside the loop, and every
©2026 HuaShu. Personal use of this material is permitted. This is an independent reformatting of the author’s open “Orange Book” guide Loop Engineering: Stop Asking Me What It Is
(v260615) into a conference-style document. Originally released June 2026; freely available at huasheng.ai/orange-books.



% %%PAGE 2

TABLE I
The Four-Layer Stack
Layer What it minds Core question
Prompt
eng.
Writing one
good prompt
What should I tell the
model
Context
eng.
What goes in the
window now
What to retrieve, sum-
marize, clear out
Harness
eng.
Arming a single
run
Which tools, which ac-
tions, what counts as
done
Loop eng. Scheduling on
the harness
How to make it run it-
self over and over
Loop engineering
make it run itself, over and over
Harness engineering
arm a single run: tools, actions, “done”
Context engineering
what goes in the window right now
Prompt engineering
the words you write for the model
scope grows
one
floor
up
Fig. 1. The four-layer stack. Each layer minds something larger than
the one below; loop engineering automates the “waiting for you” that
the harness leaves behind.
tick must come from the human. In the new world, one designs
something that ticks on its own—it runs on a timer, spawns
helpers to do the work, and feeds its own results back to itself.
As Osmani puts it, loop engineering sits one floor above the
harness: the harness below arms a single agent run; the loop
above makes it run itself over and over. The change is one of
identity, from the person who operates the agent to the person
who schedules it. Value moves from “knowing how to direct” to
“knowing how to build loops, and how to put a check inside the
loop that can say no”—the last part being the hardest, as later
sections explain.
II. F rom Prompt to Context to Loop
The “XX Engineering” terms are not replacements for one
another; they stack, each minding something larger. Table I lays
out the four layers.
Each layer up, the unit of concern grows one size: from one
sentence, to one window, to one run, and finally to a loop that
runs itself. Fig. 1 shows the four layers as a stack, with the loop
sitting one floor above the harness.
A. The Four Layers
Prompt is the bottom layer and the best known. It minds what
to tell the model: wording, examples, role, tone. Its boundary
is one exchange. The trouble is that it assumes the human is
present every time to hand the prompt in.
Context raises the question from “what do I say” to “what
should go into this window so the model can crack the problem.”
It minds the model’s entire field of view—what to retrieve, how
to summarize, what stale information to clear. A window full of
noise wastes even the best prompt.
Harness minds what an agent needs to carry for a single run:
which tools, which actions, when to load context, how to recover
from failure, what state counts as done. It arms one run; it does
not make that run repeat.
Loop automates the “waiting for you” away. With the first
three layers in place, an agent runs cleanly once but then stops.
The loop fits it with a timer so it wakes on schedule, spawns
sub-agents for parallel work, and feeds its own output back as
the next round’s input.
B. What One Floor Up Actually Adds
Three verbs separate harness from loop. Runs on a timer: the
loop wakes on schedule with no button-press. Spawns helpers: a
turning loop splits off sub-agents—one drafts a change, another
does nothing but pick it apart in review. Feeds itself: what the
loop produces becomes its own input next round; yesterday’s
findings are written to a file, and this morning it reads that file
and carries on. This memory across conversations is why it is a
loop rather than a one-off task run many times.
The fine-grained split matters because each layer fails dif-
ferently, and the check that can say “no” must be installed in a
different place. A bad prompt is caught on the spot; bad context
shows up in a wrong answer. But at the loop layer the system
runs while one sleeps, changes code one never looked at, and
feeds its own errors into the next round—and the mistake may
go undiscovered for days. The higher the layer, the farther one is
from the scene, and the longer mistakes pile up. That is precisely
why the real difficulty of loop engineering is never building the
loop, but putting something inside it that can stop it.
C. Each Layer’s Failure Has a Different Blast Radius
Consider the same underlying bug—an agent that misreads
what a function returns—as it manifests at each layer. At the
prompt layer, the misreading produces one wrong answer in
one exchange; the human sees it immediately and rewrites the
prompt. At the context layer, the misreading comes from stale
documentation loaded into the window; the human notices the
answer is confidently wrong and clears the context. At the har-
ness layer, the agent acts on the misreading once—perhaps it
edits a file—but the run ends, the diff is visible, and the human
reviews it before anything ships. At the loop layer, the same mis-
reading is written into the state file, read back the next morning
as established fact, and built upon across many turns. By the
time anyone looks, the wrong assumption is load-bearing.
This is the single most important intuition in loop engineering:
the cost of a mistake scales with the number of turns it survives
before someone catches it, and a loop is, by construction, a
machine for maximizing the number of turns. Everything in the
later sections—the evaluator, the human checkpoint, the budget
caps—exists to shorten the distance between a mistake and its
discovery.
2



% %%PAGE 3

Dis-
covery
Hand-
off
Verifi-
cation
Persis-
tence
Sched-
uling
one
turn
the “say no”
Fig. 2. The five moves of one turn. Scheduling closes the cycle—it
feeds an unfinished turn into the next day’s run. Verification is the
move that can say “no.”
III. T he Five Moves of One Loop
The word “loop” is easy to misread as idle spinning. Each
turn does something concrete: it finds work worth doing, hands
it to an agent, verifies whether the result is right, saves state,
then decides the next step. Drop any one of the five moves—
discovery, handoff, verification, persistence, scheduling—and
the loop will not turn, or will turn in place. Fig. 2 traces the five
as one turn that feeds the next.
A. A Concrete Example
Osmani built himself a morning triage loop. In the morning an
automation kicks off on its own. A triage skill reads yesterday’s
failing CI tests, the still-open issues, and recent commits, and
writes its results into a markdown file or a Linear board. For
each finding worth acting on it opens an isolated worktree; one
sub-agent drafts the fix, a second reviews it against the project’s
skills and tests. A connector automatically opens the pull request
and updates the ticket. Anything it cannot handle goes to an
inbox to wait for a human, and a state file survives so the next
day picks up where this one left off. No step needs a hand, yet it
stops to wait for a human exactly where it should.
B. The Moves
Discovery figures out what this turn should do. In the ex-
ample, the triage skill reads CI failures, open issues, and recent
commits. The key is letting the agent find its own work rather
than being handed a list. Crucially, the automation triggers a
skill—knowledge made permanent—rather than a wall of in-
structions pasted into a cron job nobody will update. Discovery
sets the ceiling on the whole loop’s quality: surface work of no
value and the other four moves are done beautifully in service of
nothing.
Handoff moves the task from the scheduling system into the
hands of the agent that does the work. Each finding worth doing
gets its own isolated git worktree, so multiple agents change
code in separate directories without stepping on each other. The
cleaner each task is cut, the easier verification and merging are
later.
Verification is the easiest move to cut corners on and the one
TABLE II
The Five Moves, Mapped to the Triage Loop
Move What it does In the triage loop
Discovery Find this turn’s
work on its own
Skill reads CI / is-
sues / commits
Handoff Hand the task o ff,
isolated
Each finding opens
a worktree
Verification Swap in another
agent to say no
Second sub-agent
reviews vs. tests
Persistence Write state outside
the conversation
PR + inbox + state
file
Scheduling Make it turn round
after round
Morning automa-
tion runs on its own
least affordable to skip. After the first sub-agent drafts the fix, a
second sub-agent reviews it—different instructions, sometimes a
different model. The agent that wrote the code grades its own
homework too softly; a dedicated hole-picker catches what the
first talked itself into letting through. This is the “thing that can
say no.” A loop without a real check is just an agent nodding at
itself.
Persistence lands the result somewhere that survives the con-
versation: a PR and updated ticket via a connector, an inbox for
what cannot be handled, and a state file recording progress. A
loop’s memory cannot live only in the context window; what is
written to markdown or a board does not forget.
Scheduling is what makes one turn into a loop. The triage
runs automatically each morning, and the state file lets unfinished
findings carry to the next day, which picks up on its own. As
Osmani puts it, automations are what make a loop an actual loop
and not just one run you did once. Table II summarizes the five.
IV . Six Parts: What a Loop Is Built From
If moves describe what happens in one turn, parts describe
what must be in hand for it to turn at all. The two line up:
discovery runs on skills, handoff on worktrees, verification on
sub-agents, persistence on memory, scheduling on automations.
Automations get the loop moving on their own, hanging off
a schedule or trigger. Without scheduling, what one holds is a
single run, not a loop. The automation should trigger a named
skill, not a wall of instructions in a cron job. Scheduling comes
in more than one form—local (machine must stay on) and cloud
(turns even with the machine off).
Worktrees are a built-in git mechanism for multiple inde-
pendent working directories in one repo. Their value scales
with parallelism: two agents writing the same file at once is the
same headache as two engineers committing to the same lines.
Worktrees turn parallelism from “runs but messy” into “runs and
clean.”
Skills make project knowledge permanent in a single file
(SKILL.md), so the agent need not re-derive context every turn.
Osmani names the cost they pay off as intent debt: the price of
explaining “what this project is, what the rules are, where the
traps are” over and over. A skill can be reused and maintained; a
wall of prompt cannot.
3



% %%PAGE 4

TABLE III
Six Parts Mapped to Five Moves
Part What it is Maps to move
Automations Runs off a schedule
/ trigger
Scheduling
Worktrees Isolated dirs for par-
allel agents
Handoff
Skills Permanent knowl-
edge; pays off intent
debt
Discovery
Connectors MCP hookup to ex-
ternal systems
Persistence / Dis-
covery
Sub-
agents
Generator separated
from judge
Verification
Memory Persistent state on
disk
Persistence
Connectors (built on MCP, the Model Context Protocol) hook
the loop to the outside world—the issue tracker, the database, a
staging API, Slack. A loop that can only see the filesystem is a
tiny loop. Connectors decide the loop’s radius of vision, and a
connector written for one tool can often be dropped onto another
and used as is.
Sub-agents split the one that writes from the one that judges.
When one agent is both player and referee, the referee plays
favorites. Counterintuitively, tuning an independent judge to be
picky is far easier than getting the generator to be hard on its
own work—which is why a loop keeps an extra agent rather than
letting one agent audit itself.
Memory is persistent state living outside a single
conversation—a markdown file or a board. The moment a con-
text window is cleared, the agent remembers nothing; for a loop
to pick up today where it left off yesterday, memory must land
on disk. The agent forgets, the repo does not. Memory is not
context: context is what the agent sees this round and is flushed
on refresh; memory persists across rounds and days. Table III
maps the six parts to the five moves.
With all six in place a loop has a skeleton: automation makes
it move, worktrees keep it from fighting itself, skills keep it
from redoing work, connectors let it see outside, sub-agents let
it correct itself, and memory lets it remember. But building it
is only the start: the same set of parts, built by two people, can
come out completely opposite.
V . Generator and Ev aluator
The hardest part of a loop is not getting the agent to run,
but putting something inside that can say “no”—and the agent
writing the code is the one least likely to say it.
A. It Always Praises Itself
Ask an agent to grade what it just produced and it tends to
praise it confidently, even when a human can plainly see the
quality is mediocre, as Anthropic engineer Prithvi Rajasekaran
observed while building long-running applications. This is not a
smarts problem; it is grading one’s own homework. The context
Generator
writes the code
Evaluator
different model;
assumes broken
draft
reject + reasons
acts via MCP: click, screenshot, run tests
maker–checker
Fig. 3. Generator and evaluator as separate agents. The evaluator
carries none of the generator’s self-persuasion, defaults to doubt, and
judges behavior by acting rather than just reading.
in which the code was written is already stuffed with the reasons
it was written that way, so when the agent looks at its own output
it does not see the result—it sees the chain of self-persuasion
that led there. Inside a loop the flaw is amplified: if every “is
this good enough” is decided by the agent that just wrote it, each
round it nods at itself, and the longer it runs the further it drifts
from real quality.
B. Tune a Skeptic, Don’t Fix a Modest Author
Making the generator more self-critical works poorly. Ra-
jasekaran found that tuning a standalone evaluator to be skeptical
is far more tractable than making a generator critical of its own
work. The difference is structural, not a matter of wording: one
cannot ask an author to step outside its own perspective, but one
can swap in another agent with entirely di fferent instructions
that looks at the code from scratch, carrying none of the self-
persuasion. The idea is borrowed from generative adversarial
networks (GANs)—one network builds, one picks faults—ported
to a generator that writes and an evaluator that reviews. Fig. 3
shows the loop this forms.
C. The Evaluator Should Act, Not Just Read
Swapping agents is not enough. If the evaluator only reads
code it judges “does this look right,” not “does it run right.” On
frontend tasks Rajasekaran hooked the evaluator to Playwright
MCP so it could open the page, click buttons, take screenshots,
and inspect the DOM like a QA engineer. That shifts the basis for
judgment from “this JSX looks fine” to “I clicked the button, the
page navigated, here is the screenshot.” Swapping the underlying
model too helps: the same model with new instructions often
keeps its blind spots. A common community calibration tells the
evaluator to assume the code is broken until proven otherwise—
the default stance should be doubt, not trust.
D. In a Product: /goal on the Stop Condition
Claude Code turns this structure into a primitive with /goal:
give an agent a condition and let it run until the condition is met.
A representative evaluator setup and stop condition look like the
following.
\# Evaluator agent (.claude/agents/reviewer.md)
ROLE: Adversarial code reviewer.
ASSUME: this code is BROKEN until proven otherwise.
DO NOT praise. Find what fails.
CHECK, in order:
4



% %%PAGE 5

1. Does it run? (execute, don’t read)
2. Tests: run them, paste real output.
3. Edge cases the author skipped.
4. Does behavior match the ticket?
USE Playwright MCP: open the page, click,
screenshot, inspect the DOM. Judge behavior,
not intent.
VERDICT: PASS only if every check holds.
Otherwise REJECT + list each reason.
\# Stop condition, judged by a fresh small model
/goal all tests in test/auth pass and the lint
step is clean
Crucially, after each turn a small fast model checks whether the
condition holds; if not, another turn runs instead of returning
control. Completion is decided by a fresh model, not the one
doing the work. This is the maker–checker principle—decades
old in banking, where the person entering a large transfer and the
person reviewing it must differ—applied to the stop condition.
(Codex reaches the same capability through automations plus
agent configuration; one should not confuse /goal with /loop,
which merely reruns on an interval.)
A loop’s floor is its evaluator. The generator’s level decides
what a loop can produce; the evaluator’s level decides what it
will not produce. Separate generation from judgment structurally,
tune the evaluator into a skeptic, make it verify by acting, and
hand the final say to a fresh model—those four steps are what it
takes to grow a loop’s ability to say “no.”
VI. F ive Ways a Loop Goes Wrong
Before turning to loops that work, it is worth cataloguing the
ways they fail, because the failures are more instructive than
the successes and far more common. Each anti-pattern below
corresponds to one of the five moves being skipped or done
badly.
A. The Nodding Loop (verification skipped)
The most common failure. The loop runs, the agent writes
code, and the same agent declares it good. With no independent
check, every turn produces self-approved output, and the loop
accumulates plausible-looking mistakes at machine speed. The
symptom is a loop that has never once said “no” to itself across
hundreds of turns—a statistical impossibility for any real work-
load, and therefore proof that no real check exists. The fix is the
generator/evaluator split of the previous section.
B. The Amnesiac Loop (persistence skipped)
The loop discovers good work, does it, and then forgets it
happened, because the result lived only in a context window that
was flushed. The next turn rediscovers the same work, or worse,
redoes it and conflicts with the first attempt. The symptom is a
loop that makes no cumulative progress: each morning it starts
from the same place. The fix is a state file on disk—the agent
forgets, the repo does not.
Discovery
Handoff
Verification
Persistence
Scheduling
Blind loop
Tangled loop
Nodding loop
Amnesiac loop
Manual loop
skip
skip
skip
skip
skip
Fig. 4. Each anti-pattern is one move skipped. The five failures map
one-to-one onto the five moves of a single turn.
C. The Manual Loop (scheduling skipped)
A loop with four good moves but no automation is not a loop;
it is a script the human runs by hand and then forgets to run
again. It works impressively the day it is built and silently stops
the day attention wanders. The symptom is a loop whose last
run was the day it was demoed. The fix is a real trigger—a timer
or an event—that does not depend on the human remembering.
D. The Blind Loop (discovery skipped)
The human still hands the loop its work each morning—“fix
these three bugs”—so the loop has automated the doing but not
the finding. This saves less than it appears to, because choosing
what to work on is often the expensive part. The symptom is
a human who is still spending their morning deciding what the
loop should do. The fix is to teach discovery into a skill so the
loop surfaces its own work.
E. The Tangled Loop (hando ff skipped)
The loop runs several agents in parallel but lets them all
change the same working directory, so their edits collide and
the merge is a mess no one can untangle. The symptom shows
up only under parallelism: a single-agent loop looks fine, and
the problem appears the first morning five agents run at once.
The fix is one isolated worktree per task. Fig. 4 maps the five
anti-patterns to the moves they violate.
These five are not independent. A loop missing verification
tends also to miss persistence, because a team careless about
one check is usually careless about the others. In practice they
cluster: the disciplined loop installs all five moves, and the hasty
loop installs only discovery and handoff—the two that produce
visible output—and skips the three that produce safety. The next
section shows three loops that installed all five.
VII. L oops That Run While You Sleep: Three Real Ones
Three public cases differ wildly in scale but share one skele-
ton: a trigger presses start, a set of constraints keeps it on the
rails, and a human checkpoint sits at the end. “Running while
you sleep” was never about how strong the model is—it is about
how solid that skeleton is.
A. One Engineer’s Morning
Osmani’s triage loop, broken down in Section III, runs au-
tomatically every morning. The one detail worth re-flagging:
the automation invokes a skill, not a giant block of instructions
5



% %%PAGE 6

Human trigger — @bot in Slack / emoji reaction
Deterministic orchestrator — scan links, pull Jira,
Sourcegraph + MCP assemble context (not the LLM)
LLM agent — writes code with materials in hand
Hard-coded gate — linter runs; agent cannot skip
LLM agent — fixes the lint
Hard-coded step — git commit
Human review — 1,300 PRs / week
Fig. 5. Stripe’s Minions pipeline. Deterministic gates (blue) and LLM
steps (green) interlock; anything rule-bound is kept out of the
probabilistic model. Reliability comes from the constraints, not model
size.
glued into a schedule that nobody will ever update. This is what
a loop one person can run looks like—one person, one machine,
the grunt work done each morning.
B. Stripe’s Minions: 1,300 PRs a Week
For enterprise scale, the case to study is Stripe’s Minions:
more than 1,300 pull requests merged a week, not one line
written by hand, as described by Stripe engineer Steve Kaliski on
the How I AI podcast. The trigger is light—@ the bot in Slack,
or add an emoji reaction. What makes it reliable is the stretch
before the model wakes up: a deterministic orchestrator first
assembles context, scanning links, pulling Jira, finding docs, and
using Sourcegraph plus MCP to locate relevant code. Letting
the LLM find its own context is the least controllable part, so
that work—whose rules can be hard-coded—is taken out of the
model’s hands. Anything deterministic logic can solve never
goes to a probabilistic model; where one draws that line decides
whether the loop is reliable.
The most counterintuitive point: Minions is not built on a
stronger model. It is a fork of the open-source tool Goose, and
its core claim is that reliability comes from the quality of the
constraints, not the size of the model. Its architecture inter-
leaves deterministic gates and creative LLM steps, as sketched
in Fig. 5—the agent writes code, a hard-coded pipeline runs
the linter and the agent cannot skip it, the agent fixes the lint,
and a hard-coded step runs the commit. The sandbox is Devbox
on EC2, run on a “cattle not pets” basis: each environment is
swapped out at will, so a thousand-plus agents run at once with-
out stepping on each other. Notably, those 1,300 PRs are still
reviewed by humans—the human did not leave, but changed
desks, from writing to reviewing.
C. What “While You Sleep” Actually Relies On
Local /loop and desktop scheduled tasks need the machine
on; turn it off and the loop stops. To run with the machine off,
the right answer is Cloud Routines or GitHub Actions schedule
triggers. Table IV contrasts the options. Want it frequent and
able to see local files? Use local /loop, at the cost of keeping
TABLE IV
Scheduling Options Compared
Cloud Desktop /loop
Where it runs cloud machine machine
Machine on? no yes yes
Session open? no no yes
Min. interval 1 h 1 min 1 min
See local files? no yes yes
the machine on. Want it to run untethered from local state? Go
to the cloud, at the cost of a one-hour minimum interval and a
clean clone each time. No single scheduler does it all.
A caution on widely circulated numbers: claims such as
“around 90\% of Claude Code is written by itself” or large migra-
tion speedups are mostly secondhand summaries and should be
treated as rough reference. The three cases here trace to firsthand
sources, which holds up better than one impressive-sounding
figure.
D. Choosing a Scheduler in Practice
The choice between local and cloud scheduling is not a matter
of taste; it follows mechanically from one question: is the loop’s
work glued to the local machine, or can it leave? Two concrete
scenarios make the rule clear. Suppose a loop must check a
local development server every minute—that work can only run
locally, because the cloud cannot see a process on one’s laptop
and the cloud interval cannot drop below an hour. Now flip it:
suppose a loop should scan the repository’s open issues at three
in the morning and open pull requests where warranted—that
work should never be tied to a laptop at all, because laptops get
their lids closed, lose power, and get carried out the door. For
the second, a cloud schedule or a CI schedule trigger is the right
answer, running on a machine that stays awake while the human
sleeps.
The distortion to avoid is treating local rerun as the whole of
“running while you sleep.” A local rerun means “run a few extra
rounds while I am here”; cloud scheduling means “run even
when I am not.” These are different capabilities, and conflating
them is how people end up disappointed when they close the
lid and the loop they thought was autonomous quietly stops.
The honest framing is that local scheduling buys frequency and
access to local files at the cost of requiring the machine to stay
on, while cloud scheduling buys true autonomy at the cost of a
coarser interval and a fresh clone each run. No single scheduler
does it all, and a mature loop often uses both—local for the tight
inner checks, cloud for the overnight sweep.
E. The Same Capability, Two Toolchains
The commands in this note are Claude Code’s, but the capa-
bilities are not specific to it. Codex offers the same five organs
under different names, and a connector written for one side can
often be moved to the other and used as is. Table V lines them
up so the reader does not pin a command name on the wrong
door. The lesson is that loop engineering is a set of capabilities,
not a product: scheduling, run-until-condition, parallel isolation,
6



% %%PAGE 7

TABLE V
The Same Capability Across Two Toolchains
Capability Claude Code Codex
Scheduling /loop worker Automations tab
Run until met /goal automation
rerun + judge
Parallel isolation --worktree background
worktree
Sub-agents .claude/agents/ .codex/agents/
External conn. MCP + plugins MCP connector
Explicit skill SKILL.md \$skill-name
Machine-off run Cloud Routines cloud (planned)
sub-agents, external connection, and explicit skill invocation.
Whichever toolchain a team uses, the question to ask is whether
all six are present, not which brand of command provides them.
VIII. T he Costs: Four Tabs That Don’t Clear Themselves
A loop that runs itself is, at the same time, a loop that makes
mistakes by itself. The more cheerfully it runs, the more quietly
it errs. Four costs accrue, none of which sounds an alarm while
the loop is running.
Verification debt. Every PR opened and merged saves time,
but the saved time turns into unverified output waiting to be paid
back. The problem hides where tests do not cover, in the gap
between “runs” and “right,” accumulating until some shipping
morning when it blows up at once. The guard is an independent
evaluator—a different agent from the one doing the work.
Comprehension rot. The faster the loop ships code one did
not write, the bigger the gap between what exists and what one
actually understands. Reading code is more boring than writing
it, and the loop has taken the writing; the codebase grows while
the map in one’s head stalls. It sounds no alarm until a bug
burrows into a corner never read. The guard is to read the loop’s
output regularly and force oneself to explain a few changes; an
inability to explain is a map needing an update.
Cognitive surrender. When the loop runs itself it is tempting
to stop having an opinion and just take whatever it hands back.
This is the attitude version of the first two: not “no time” but “no
longer want to bother.” The more reliable the loop, the easier it
is to outsource judgment. The guard is one line—the loop can
execute, but it cannot decide. One must at least remain capable
of saying “this is wrong.”
Token blowout. The only cost that hits the bill directly, and
hard to estimate in advance: the loop hatches helpers, retries,
and runs round after round, so one bug can spin idle all night
and produce an unfamiliar bill rather than fixed code. The guard
is hard caps set before shipping—per-run budget, daily budget,
max retries—so an idle bug cannot burn an entire night’s quota.
The four share one trait: silence while the loop runs. The
most fascinating thing about loop engineering is that it lets one
person do a team’s work; the most dangerous thing is the same
spot, because a team argues with itself and one person plus a pile
of loops easily becomes an echo chamber where no one argues.
Verification
debt
Comprehension
rot
Cognitive
surrender
Token
blowout
each feeds
the next
Fig. 6. The four costs reinforce one another. Unverified output erodes
understanding, which invites surrender, which lets the loop run longer
and spend more, which produces more unverified output.
A. A Worked Example of Compounding Debt
Consider a loop that opens twenty pull requests overnight, all
with green tests. On the surface this is a triumph. But suppose
three of the twenty contain a subtle error the tests do not cover.
With no independent evaluator, those three merge—that is veri-
fication debt. Because the human merged twenty PRs without
reading them, their mental model of the codebase now lags by
twenty changes—that is comprehension rot. Because the loop
ran so smoothly, the human stops reading the next morning’s
batch entirely—that is cognitive surrender. And because the loop
spawned helpers and retried freely all night, the bill is triple what
was estimated—that is token blowout. The three buried errors
now sit in a codebase the human no longer fully understands,
guarded by a human who has stopped looking, discovered even-
tually only when one of them surfaces as a production incident.
The point of the worked example is that the four costs are not
a list of independent risks; they are a single failure that wears
four faces. They reinforce one another: the more output the
loop produces unverified, the less the human understands; the
less they understand, the more they surrender; the more they
surrender, the longer the loop runs unwatched and the bigger
the bill. Fig. 6 shows the reinforcing cycle. The guard against
all four is the same: keep a human capable of saying “no,” and
install a check the human does not have to be awake to run.
IX. S tay the Engineer, Not Just the One Who Presses Go
The same loop, built by two di fferent people, can end in
opposite places—and the difference is not in the loop. As Osmani
writes, two people can build the same loop and get opposite
outcomes. One uses a loop to move faster on things already
mastered: they read the code, hold a firm sense of direction, and
the loop scales judgment they already had. Another uses the
same loop so they never have to understand again. Six months
later, one has gotten stronger and the other has become the
gatekeeper of a machine they cannot read.
A loop is not a tool whose quality is fixed by the tool. It is so
strong that it amplifies, unchanged, whatever one brings: bring
understanding and it amplifies understanding; bring laziness and
it amplifies laziness. It is a faithful multiplication sign, and what
it multiplies is the person.
The loop makes generation extremely cheap—code, plans,
7



% %%PAGE 8

PRs, fixes, almost free. What stays scarce is judgment: knowing
which plan is right, which line should be stopped, which output
runs fine but is wrong at the root. The loop can generate a
hundred options but cannot truly choose; or rather, it chooses
on “looks reasonable,” not “actually right,” and the gap between
those two is the reason an engineer exists. Loop engineering
therefore does not devalue judgment; it strips away everything
that does not require judgment and leaves judgment as all that
remains.
The loop executes the logic given to it, but does not under-
stand why one wanted to build it, what one actually wants, or
which spots one would rather watch oneself. Those boundaries—
where a manual checkpoint sits, where it runs on its own—cannot
be read out of the loop; they live only in the builder’s head and
must be written in one by one. The mindset one designs with
is the shape the loop grows into. Two loops built with a “free
myself fast” mindset and an “I still mean to be the engineer”
mindset may be ninety percent identical in code; the difference
is one or two checkpoints, and those decide whether, six months
out, one stands on top of the loop or is hollowed out by it. Os-
mani’s closing line is the one to keep: build the loop, but build
it like someone who intends to stay the engineer, not just the
person who presses go.
X. T he Economics of Judgment
The previous section made a claim worth examining on its
own terms: that loops make generation cheap and leave judgment
scarce. This is not a slogan; it is an economic observation with
consequences for how a team should organize.
A. What Becomes Abundant
When a resource becomes abundant, its price falls and the
activities organized around it reorganize. Loops make code,
plans, fixes, and pull requests abundant—a single engineer with
a well-built loop can produce the output of a small team. The
activities that used to consume an engineer’s day, the typing and
the boilerplate and the mechanical refactor, collapse toward zero
cost. A reasonable first reaction is that this makes engineers less
valuable. The opposite is true, but only for engineers who hold
onto the scarce thing.
B. What Stays Scarce
The scarce resource is the judgment that decides which of
the abundant outputs to keep. A loop can generate a hundred
candidate implementations; it cannot tell you which one is right,
only which one looks reasonable, and the gap between “looks
reasonable” and “is right” is exactly where engineering lives.
As generation approaches free, the entire value of the engineer
concentrates into this gap. The work that remains is pure judg-
ment, distilled and undiluted by the mechanical labor that used
to surround it.
This has an uncomfortable implication. An engineer whose
value was mostly in the mechanical labor—fast typing, broad
memorization of APIs, willingness to grind through boilerplate—
finds that value evaporating, because the loop does all of it for
free. An engineer whose value was in judgment finds it amplified,
because the loop executes their good decisions a hundredfold.
The same tool widens the gap between the two kinds of engineer.
It does not lift everyone equally; it multiplies whatever each
person brings.
C. The Amplifier Cuts Both Ways
Because the loop is an amplifier of judgment, a lapse in
judgment is also amplified. In the old world a bad decision cost
one hand-written stretch of wrong code, limited in blast radius
and slow enough to catch. In the new world a bad decision is
executed faithfully, in bulk, a hundred times, by a machine that
will not pause to ask whether it is right. The loop removes the
slow gear that used to bail engineers out. One can no longer
count on the process being slow enough to notice a mistake mid-
flight, because the process has no slow gear left. This raises the
stakes on the one thing the loop cannot do, and it is the reason
the discipline of staying the engineer is not optional sentiment
but operational necessity.
XI. O perational Discipline
If judgment is the scarce resource, the practical question is
how to spend it well. Three disciplines, drawn from the cases
and costs above, are worth stating as standing practice.
A. Read a Sample, Always
The defense against comprehension rot is not to read every-
thing the loop produces—that would defeat the purpose—but
to read a representative sample, every day, and to force oneself
to explain each sampled change: what it did and why it did it
that way. An inability to explain a change is a precise signal that
one’s mental map has fallen behind the codebase, and it is far
cheaper to discover this from a sampled PR on a quiet morning
than from a production incident on a bad one. The sample need
not be large; it needs to be regular and genuinely examined.
B. Cap Before You Ship
The defense against token blowout is to set hard ceilings
before the loop runs unattended for the first time, not after the
first surprising bill. A per-run budget, a daily budget, and a
maximum retry count together ensure that a single bug spinning
idle overnight cannot consume an entire quota. These numbers
are not primarily about saving money; they are circuit breakers
that convert an open-ended risk into a bounded one. A loop
without caps is a loop that has delegated its spending authority
to its own bugs.
C. Keep One Door Open
The defense against cognitive surrender is structural, not
merely attitudinal. Build at least one checkpoint into the loop
where it pauses for a human—not because the human will always
intervene, but because the existence of the pause keeps the human
in the position of being able to. The engineer who welds every
door shut, banking on never needing to go in, discovers on the
day they must that they no longer hold the key. The engineer
who leaves one door open can walk in any time to see what the
loop is doing. The two loops di ffer by a single checkpoint in
8



% %%PAGE 9

code; they di ffer enormously in who is in control six months
later.
XII. B uild Your First Loop Today
Stripe’s pipeline is the endpoint, not the starting point. A first
loop should be so small it barely looks like a system—a little
thing that checks something on a timer.
Step one: run a /loop. Available after Claude Code v2.1.72,
it reruns the same task on an interval. It is session-scoped,
recurring tasks expire after seven days, and it runs on the local
machine; turn the machine off and it stops.
/loop 5m check the deploy \# fixed: every 5 min
/loop check the deploy \# agent paces itself
/loop \# runs .claude/ loop.md
Step two: read CI and issues; triage first. Rerunning one
line is not a loop. Give it a prompt to look at three things each
morning and have it list what is worth handling. Scheduled plus
auto-discovery is loop entry level. The discovery logic should
live in a skill, not the schedule.
\# .claude/skills/morning-triage/SKILL.md
NAME: morning-triage
WHEN: invoked each morning by automation.
READ:
- CI runs that failed since yesterday
- issues opened in the last 24h
- commits merged since the last run
JUDGE: for each item, is it worth acting on?
Skip noise. Keep only actionable findings.
OUTPUT: write findings + status to
./state/triage.md (one row per finding).
Step three: add a state file. Do not leave results in the chat
window. Write every finding, and how far it has been handled,
into a markdown file (or a Linear board). The agent forgets; the
repo does not.
\# ./state/triage.md (the loop’s memory)
| finding | source | status |
|----------------|----------|----------|
| auth test flaky| CI \#4821 | fixing |
| null deref | issue 92 | PR open |
| stale dep | commit a3| inbox |
Step four: add an evaluator. The most critical step and the
easiest to skip. Claude Code’s /goal (after v2.1.139) runs until
a condition is met, with a di fferent model judging whether it
holds.
/goal all tests in test/auth pass and the lint
step is clean
Step five: add worktrees for parallelism. Use --worktree
(or -w) to open an independent worktree per background agent
so they do not step on each other.
\# one isolated worktree per finding
claude -- worktree fix/auth-test "draft the fix"
claude -- worktree fix/null-deref "draft the fix"
Of the six elements in Table VI, the first two decide whether
TABLE VI
First-Loop Checklist
Element Ask yourself
Discovery
source
What does it read on a timer? (CI /
issues / commits / inbox)
State file Which disk file holds the cross-
round memory?
Evaluator Is there an independent check that
can say “no”?
Isolation Does each parallel agent get its own
worktree?
Token cap Did you set a spending ceiling?
Who stops it if it runs off?
Human review Which step pauses for you to look,
rather than auto-ing all the way
through?
the loop can run; the last four decide whether it gets into trouble
once it does. Beginners most often ship with only the first two
built, and the result is a loop nobody watches and nobody can
stop, nodding at itself. The closing advice is simple: a first loop
is better small, but with the “no”-saying check and the human
review point fully installed.
A. A Complete First Loop, Annotated
To make the checklist concrete, the following is a minimal but
complete loop that installs all six elements. It is small enough
to read in one sitting and contains every organ a real loop needs,
only scaled down.
\# 1. SCHEDULING -- a real trigger
\# (.github/workflows/triage.yml)
on:
schedule:
- cron: ’0 6 * * *’ \# 06:00 daily, cloud
\# 2. DISCOVERY -- a skill, not a wall of text
\# invoked by the workflow:
run: claude -- skill morning-triage
\# 3. PERSISTENCE -- state on disk
\# the skill writes ./state/triage.md
\# and commits it back to the repo
\# 4. HANDOFF -- one worktree per finding
for finding in \$(parse ./state/triage.md); do
claude -- worktree "fix/\$finding" \textbackslash\{\}
--goal "tests pass and lint is clean" \textbackslash\{\}
"draft a fix for \$finding"
done
\# 5. VERIFICATION -- a fresh model judges
\# / goal’s stop check runs after each turn;
\# a second reviewer agent picks holes
\# 6. HUMAN REVIEW -- the open door
\# PRs are opened, never auto-merged;
\# anything uncertain lands in ./inbox/
Read top to bottom, the six numbered comments are the six
elements of the checklist, each realized in two or three lines.
The cron line is scheduling; the skill invocation is discovery; the
committed state file is persistence; the per-finding worktree is
handoff; the /goal stop check plus reviewer is verification; and
the never-auto-merge rule with an inbox is the human review
9



% %%PAGE 10

point. A loop with all six, even a tiny one, is a real loop. A
loop missing any of them is one of the five failures of Section VI
wearing a disguise.
B. Growing the Loop Safely
Once the minimal loop runs, the temptation is to scale it—
more findings, more parallel agents, shorter intervals. The safe
order of growth is to add parallelism last, after the checks are
proven. Increase what the loop discovers before increasing how
much it does in parallel, and prove the evaluator catches real mis-
takes before trusting it to gate many agents at once. The Stripe
case is the endpoint of this path, not the entry: its reliability
comes from years of hardening the deterministic gates, not from
starting large. A loop earns the right to run more agents by first
demonstrating it can stop a single bad one.
The rest is not in this note; it is in the terminal.
XIII. A ppendix A: An AnnotatedTriage Skill
The discovery move depends on a skill rather than a wall of
instructions, because a skill can be reused and maintained while
a pasted prompt rots in a schedule nobody updates. The fol-
lowing is a fuller version of the morning-triage skill referenced
throughout, annotated to show how each section serves a move.
\# .claude/skills/morning-triage/SKILL.md
---
name: morning-triage
trigger: invoked by daily automation
---
\#\# Read (the DISCOVERY inputs)
- CI runs failed since the last run
- issues opened in the last 24 hours
- commits merged since yesterday
- the previous ./state/triage.md
\#\# Judge (the part that sets the ceiling)
For each candidate, decide:
- is it actionable now, or noise?
- does it block a release? → priority
- is it already tracked? → skip
Keep only what is worth a worktree today.
\#\# Write (the PERSISTENCE output)
Append to ./state/triage.md:
| finding | source | priority | status |
Commit the file so tomorrow can read it.
\#\# Hand off (prepare the HANDOFF)
For each kept finding, emit a task line:
worktree=fix/<slug> goal=<stop-condition>
\#\# Stop (the boundary you keep for yourself)
Never merge. Never delete. Anything you are
less than confident about goes to ./inbox/
for a human, not into a PR.
Five of the skill’s six headings map to the five moves; the sixth,
“Stop,” is where the builder writes in the boundary the loop
cannot infer. The loop will faithfully do everything the skill says
and nothing it omits, so the “Stop” section is not boilerplate—it
is the single place where the engineer’s intent about where to
keep control is made permanent. Leave it out and the loop will
merge with confidence it has not earned.
TABLE VII
Terms Used in This Note
Term Meaning
Loop A system that discovers, does, verifies,
persists, and reschedules work without a
human in the inner cycle.
Harness The kit arming a single agent run: tools,
allowed actions, recovery, “done.”
Move One of the five steps in a single turn of a
loop.
Part One of the six components that realize
the moves.
Generator The agent that writes.
Evaluator A separate agent that judges, defaulting
to doubt and acting to verify.
Worktree A git mechanism giving each parallel
agent its own working directory.
Skill Project knowledge made permanent in a
SKILL.md file.
Connector An MCP interface linking the loop to ex-
ternal systems.
Memory Persistent state on disk, surviving any sin-
gle conversation.
Intent debt The recurring cost of re-explaining a
project, paid off by a skill.
Verification
debt
Unverified output accumulating between
“runs” and “right.”
XIV . Appendix B: Glossary
XV . Synthesis: What the Playbook Comes Down To
Across nine sections the argument has one spine. Loop engi-
neering is the fourth layer of a stack that began with the prompt
and climbed through context and the harness, and what dis-
tinguishes it from the three below is that it removes the hu-
man from the position of doing the work. A single turn of a
loop is five moves—discovery, handoff, verification, persistence,
scheduling—realized by six parts, and the failures of a loop are
simply those moves skipped. The hardest of the moves is verifi-
cation, because an agent grading its own work praises it, and the
reliable remedy is structural: a separate evaluator that defaults to
doubt, acts rather than reads, and is judged by a fresh model on
an explicit stop condition. Loops already run in practice, from
one engineer’s morning to an enterprise merging over a thousand
machine-written pull requests a week, and what makes them reli-
able is the quality of their constraints, not the size of their model.
They run up four silent debts—verification debt, comprehension
rot, cognitive surrender, and token blowout—that reinforce one
another and come due all at once. And because the loop is an
amplifier of whatever the builder brings, the same loop built by
two people yields opposite outcomes, separated by one or two
checkpoints that decide who is in control later.
The single sentence to carry away is the one the field con-
verged on in a week: stop prompting the agent, and design the
system that prompts it—but design it like someone who intends
to stay the engineer, not just the one who presses go. Everything
10



% %%PAGE 11

technical in this note serves that one posture. The evaluator, the
state file, the budget cap, the open door: each is a way of keeping
a human capable of saying “no” to a machine built to say “yes”
at speed. A loop is the most powerful tool in this generation of
software practice precisely because it is the most faithful multi-
plier of its builder, and a faithful multiplier is exactly as valuable,
or as dangerous, as the judgment fed into it.
A. Field Notes for the First Month
For a team adopting loops, a few practical observations recur
often enough to be worth stating plainly. First, the loop that sur-
vives is the small one that earned trust, not the ambitious one that
demanded it; start with a single finding handled end to end and
widen only after the checks have caught real mistakes. Second,
the evaluator is where the engineering effort belongs—a strong
generator with a weak judge produces confident garbage, while
a modest generator with a sharp judge produces slow, reliable
progress, and the second is what compounds. Third, the human
review point is not a temporary scaffold to be removed once the
loop is trusted; it is the permanent feature that keeps the loop
trustworthy, and the day it is removed is the day comprehen-
sion rot begins in earnest. Fourth, budget caps should be set on
the assumption that something will spin idle overnight, because
eventually something will, and the cap is the difference between
a curiosity in the logs and a line item on an invoice.
None of these is surprising in isolation. What surprises teams
is how quickly the pleasant experience of a loop that “just works”
erodes the disciplines that made it work, and how the erosion is
invisible until the morning it is not. The playbook, in the end, is
less about building loops—that part is genuinely easy now—and
more about remaining the kind of engineer who can still answer,
on any given morning, whether the thing the loop just did was
actually right.
References
[1] A. Osmani, “Loop Engineering,” personal blog and Substack, Jun. 2026.
[2] P. Steinberger, post on designing loops that prompt coding agents, social
media, Jun. 2026.
[3] B. Cherny, public remarks on writing loops that prompt Claude, Anthropic,
Jun. 2026.
[4] P. Rajasekaran, “Building long-running agentic applications: the genera-
tor/evaluator pattern,” Anthropic engineering blog, 2026.
[5] S. Kaliski, “Stripe’s Minions: 1,300 PRs a week,” How I AI podcast, 2026.
[6] “Model Context Protocol (MCP) specification,” open standard, 2025–2026.
[7] “Goose: an open-source agent framework,” project documentation, 2025–
2026.
[8] “Claude Code documentation: /loop, /goal, worktrees, skills, automa-
tions,” Anthropic, 2026.
[9] HuaShu, Loop Engineering: Stop Asking Me What It Is , Orange Books,
v260615, Jun. 2026.
Acknowledgment. This is an independent conference-style synthesis built on
the framework of HuaShu’s open guide Loop Engineering: Stop Asking Me
What It Is (Orange Books, v260615, June 2026). The framework and quoted
formulations are due to Addy Osmani; the generator/evaluator findings to Prithvi
Rajasekaran (Anthropic); and the enterprise case to Steve Kaliski (Stripe). All
product details may change; refer to each tool’s official documentation.
11
\end{document}
